\documentclass[border=5mm]{standalone}
% ==============================================================================
% SETUP.TEX - CONFIGURACIÓN GLOBAL DEL PROYECTO
% ==============================================================================
% Este archivo centraliza todos los paquetes y estilos.
% Debe incluirse en el preámbulo del libro principal y de los archivos standalone.
% \PassOptionsToPackage{gray}{xcolor}
% --- 1. CODIFICACIÓN E IDIOMA ---
% fix-cm habilita las fuentes Computer Modern en tamaños arbitrarios (por debajo
% de 5 pt, entre otros). Debe cargarse antes que fontenc.
\usepackage{fix-cm}
\usepackage[utf8]{inputenc}
\usepackage[T1]{fontenc}
\usepackage[spanish, es-tabla]{babel}  
\usepackage{csquotes} % Recomendado para babel y citas

\usepackage{import}
% mode=image: Usa el PDF pre-compilado, nunca intenta recompilar.
% Debes compilar cada figura manualmente antes de compilar el libro.
\usepackage[mode=image, subpreambles=false]{standalone}

% --- 2. MATEMÁTICAS Y CIENCIAS ---
\usepackage{amsmath, amssymb, amsfonts, mathtools}
\usepackage{enumitem}

% LaTeX trae \DeclareMathSizes{5}{5}{5}{5}, así que a 5 pt (\tiny) los subíndices
% salen del mismo tamaño que la base: en las etiquetas de figuras el M de
% $\omega_M$ queda desproporcionado. Se restablece la proporción habitual.
\DeclareMathSizes{5}{5}{3.7}{3}

% --- 3. DISEÑO Y GEOMETRÍA --- 
% Unificamos los márgenes para todo el documento
\usepackage[left=2.5cm,right=2.5cm,top=3cm,bottom=3cm]{geometry}
\makeatletter
\ifstandalone % Evita que geometry dañe el recorte de las figuras bajándolas 3cm
  \@ifclasswith{standalone}{crop=false}{
    % Es un capítulo/sección: Dejamos que geometry actúe.
  }{
    % Es una figura: Neutralizamos geometry para que no las recorte mal.
    \geometry{pass}
  }
\fi
\makeatother
\usepackage{parskip} % Espaciado entre párrafos sin sangría
\usepackage{float}   % Para usar [H] en figura

% Ajustes de flotantes para evitar espacios verticales exagerados
\renewcommand{\topfraction}{0.95}      % Permite que las figuras ocupen hasta el 95% superior
\renewcommand{\bottomfraction}{0.95}   % Permite que las figuras ocupen hasta el 95% inferior
\renewcommand{\textfraction}{0.05}     % Permite hojas con tan solo un 5% de texto
\renewcommand{\floatpagefraction}{0.8} % Requiere que una página de solo figuras esté 80% llena
\setcounter{topnumber}{4}
\setcounter{bottomnumber}{4}
\setcounter{totalnumber}{10}

% --- 4. GRÁFICOS Y TABLAS ---
\usepackage{graphicx}
\usepackage{booktabs} % Tablas profesionales
\usepackage{caption}  % Control de captions y \captionof
\usepackage{tikz}
\usetikzlibrary{arrows.meta, calc, angles, quotes, babel, backgrounds}
\usepackage{pgfplots}
\usepgfplotslibrary{groupplots}
\usepgfplotslibrary{external}
\usepgfplotslibrary{patchplots}
\pgfplotsset{compat=1.18}
\usepackage[american, straightvoltages]{circuitikz}

% --- 5. COLORES Y CAJAS ---

\usepackage[table]{xcolor}
\usepackage{tcolorbox}

% Definición de colores corporativos/estilo
\definecolor{utnblue}{RGB}{0, 51, 102}
\definecolor{codegreen}{rgb}{0,0.6,0}
\definecolor{codegray}{rgb}{0.5,0.5,0.5}
\definecolor{codepurple}{rgb}{0.58,0,0.82}
\definecolor{backcolour}{rgb}{0.96,0.96,0.96}

% --- 6. ENTORNOS PERSONALIZADOS (CAJAS) ---

% Caja "Resultado" (Azul Corporativo) — Definiciones, fórmulas clave, resultados importantes.
% Uso: \begin{resultado}[Fórmula de Euler] ... \end{resultado}
\newtcolorbox{resultado}[1][]{
    colback=utnblue!5!white,
    colframe=utnblue,
    title=\textbf{#1},
    fonttitle=\bfseries,
    sharp corners,
    boxrule=0.5mm
}

% Caja "Nota" (Gris) — Advertencias, aplicaciones, contexto secundario.
% Uso: \begin{nota}[Cuidado con la calculadora] ... \end{nota}
% Uso: \begin{nota} ... \end{nota}  → título por defecto "Nota"
\newtcolorbox{nota}[1][Nota]{
    colback=codegray!5!white,
    colframe=codegray,
    title=\textbf{#1},
    fonttitle=\bfseries,
    sharp corners,
    boxrule=0.5mm
}

% Caja inline para resaltar un valor y enlazarlo con su otra aparición en una tabla
% o en una cuenta: mismo color, mismo valor.
% Uso: \marcavalor{$4$}  o  \marcavalor[codegreen]{$4$}
\newtcbox{\marcavalor}[1][utnblue]{
    on line,
    boxsep=1pt, left=2pt, right=2pt, top=1pt, bottom=1pt,
    colback=#1!10!white,
    colframe=#1,
    boxrule=0.4pt,
    arc=2pt
}

% --- 7. CÓDIGO FUENTE (LISTINGS) ---
\usepackage{listings}

\lstdefinestyle{miestilo}{
    backgroundcolor=\color{backcolour},   
    commentstyle=\color{codegreen},
    keywordstyle=\color{magenta},
    numberstyle=\tiny\color{codegray},
    stringstyle=\color{codepurple},
    basicstyle=\ttfamily\footnotesize,
    breakatwhitespace=false,         
    breaklines=true,                 
    captionpos=b,                    
    keepspaces=true,                 
    numbers=left,                    
    numbersep=5pt,                  
    showspaces=false,                
    showstringspaces=false,
    showtabs=false,                  
    tabsize=2,
    frame=single,                    
    rulecolor=\color{codegray}
}

\lstset{style=miestilo}
% Soporte para tildes y ñ en bloques de código
\lstset{literate={ñ}{{\~n}}1 {á}{{\'a}}1 {é}{{\'e}}1 {í}{{\'i}}1 {ó}{{\'o}}1 {ú}{{\'u}}1}

% Operadores matemáticos personalizados

% Vector en sentido discreto (lista finita de coordenadas): negrita + flecha.
% Uso: \vect{v}, \vect{e}_k, \vect{0}.
\newcommand{\vect}[1]{\vec{\mathbf{#1}}}

% Fecha de versión y comando para capítulos standalone
\providecommand{\verdate}{\the\year.\ifnum\month<10 0\fi\the\month.\ifnum\day<10 0\fi\the\day}
\newcommand{\chapterversion}{%
    \vspace{-1.5cm}\begin{flushright}\small\textit{\color{codegray}Versión~\verdate}\end{flushright}\vspace{0.5cm}%
}

% --- 8. HIPERVÍNCULOS (DEBE SER EL ÚLTIMO PAQUETE) ---
\usepackage[colorlinks=true, linkcolor=utnblue, urlcolor=utnblue, citecolor=utnblue]{hyperref}

% --- 9. REFERENCIAS CRUZADAS ENTRE CAPÍTULOS STANDALONE ---
% Cuando se compila un capítulo standalone, xr-hyper lee los .aux de los
% otros capítulos (si existen) para resolver \ref{} a labels externos.
% En la compilación del libro completo este bloque no se activa.
\ifstandalone
  \usepackage{xr-hyper}
  \makeatletter
  \newcommand{\xrchap}[1]{%
    \edef\xrc@self{\detokenize{Capitulo#1}}%
    \edef\xrc@job{\jobname}%
    \ifx\xrc@self\xrc@job\else
      \IfFileExists{../#1capitulo/Capitulo#1.aux}%
        {\externaldocument{../#1capitulo/Capitulo#1}}{}%
    \fi
  }
  \makeatother
  \xrchap{01}\xrchap{02}\xrchap{03}\xrchap{04}
  \xrchap{05}\xrchap{06}\xrchap{07}\xrchap{08}
  \xrchap{09}\xrchap{10}\xrchap{11}\xrchap{12}
  \xrchap{13}
\fi
% \PassOptionsToPackage{gray}{xcolor}

% fig_margen_muestreo
% Por qué en la práctica ws tiene que superar 2 wr y no 2 wM.
% El filtro antialiasing real no corta de manera abrupta: deja pasar la banda útil
% hasta wM, decae a lo largo de una franja de transición y en wr ya se anuló.
%   X(w) = 1                para |w| < wM
%        = coseno alzado    para wM <= |w| <= wr   (1 en wM, 0 en wr)
%        = 0                para |w| > wr
% Valores: wM = 1,8   wr = 3,0   ws = 7,0  =>  2 wr = 6,0 < 7,0 (condición OK)
%          ws/2 = 3,5   ws - wr = 4,0   ws - wM = 5,2
% (a) El espectro que entrega el filtro: banda útil plana y transición (gris).
% (b) Muestreado con ws: la vecina empieza en ws - wr = 4,0, a la derecha de
%     wr = 3,0, de modo que las copias no se solapan. La regla de arriba parte el
%     eje en tres tramos: banda útil [0, wM], transición [wM, wr] y franja libre
%     [wr, ws - wr]. El rechazo del filtro está idealizado a cero; uno real deja
%     un resto por encima de wr que sí se solapa.

\begin{document}
\begin{tikzpicture}[>={Latex[length=4pt]}, font=\scriptsize, x=0.52cm, y=1.20cm,
    declare function={
        X(\x) = (abs(\x)<1.8) * 1
              + (abs(\x)>=1.8)*(abs(\x)<=3.0) * (0.5*(1+cos(180*(abs(\x)-1.8)/1.2)));
    },
    esp/.style={utnblue, thick, fill=utnblue!12},
    espcopy/.style={utnblue!70, thick, fill=utnblue!6},
    nyq/.style={gray!70, thin, dashed},
    med/.style={<->, gray!65},
    ax/.style={->, thin, black},
]
    % ===================== (a) lo que entrega el filtro =====================
    \begin{scope}[shift={(0,0)}]
        \node at (-10.4,0.62) {(a)};
        \begin{scope}
            \clip (-9.5,-0.05) rectangle (9.5,1.40);
            \draw[esp] (-3.0,0)
                -- plot[domain=-3.0:3.0, samples=300, variable=\x] (\x,{X(\x)})
                -- (3.0,0) -- cycle;
            % franja de transición, resaltada sobre el relleno
            \fill[gray!35] (1.8,0) -- plot[domain=1.8:3.0, samples=80, variable=\x]
                    (\x,{X(\x)}) -- (3.0,0) -- cycle;
            \fill[gray!35] (-1.8,0) -- plot[domain=-1.8:-3.0, samples=80, variable=\x]
                    (\x,{X(\x)}) -- (-3.0,0) -- cycle;
        \end{scope}
        \draw[ax] (-9.9,0) -- (10.1,0) node[right] {$\omega$};
        \draw[ax] (0,-0.05) -- (0,1.15);
        \node[below, font=\tiny] at (0,-0.05) {$0$};
        \draw[thin] (1.8,0.04) -- (1.8,-0.04) node[below, font=\tiny] {$\omega_M$};
        \draw[thin] (3.0,0.04) -- (3.0,-0.04) node[below, font=\tiny] {$\omega_r$};
        \node[gray!60!black, anchor=west, font=\tiny] at (4.6,0.56) {transición};
        \draw[gray!60!black, thin, -{Latex[length=3pt]}] (4.5,0.55) -- (2.6,0.34);
        \node[utnblue, anchor=south east, font=\tiny] at (-0.2,1.00) {$X(j\omega)$};
    \end{scope}

    % ===================== (b) después de muestrear =====================
    \begin{scope}[shift={(0,-2.60)}]
        \node at (-10.4,0.62) {(b)};
        \begin{scope}
            \clip (-9.5,-0.05) rectangle (9.5,1.40);
            \foreach \c in {-7,7} {
                \draw[espcopy] (\c-3.0,0)
                    -- plot[domain=\c-3.0:\c+3.0, samples=300, variable=\x] (\x,{X(\x-\c)})
                    -- (\c+3.0,0) -- cycle;
            }
            \draw[esp] (-3.0,0)
                -- plot[domain=-3.0:3.0, samples=300, variable=\x] (\x,{X(\x)})
                -- (3.0,0) -- cycle;
            % la transición de la copia central, sombreada como en (a)
            \fill[gray!35] (1.8,0) -- plot[domain=1.8:3.0, samples=80, variable=\x]
                    (\x,{X(\x)}) -- (3.0,0) -- cycle;
            \fill[gray!35] (-1.8,0) -- plot[domain=-1.8:-3.0, samples=80, variable=\x]
                    (\x,{X(\x)}) -- (-3.0,0) -- cycle;
        \end{scope}
        \draw[ax] (-9.9,0) -- (10.1,0) node[right] {$\omega$};
        \draw[ax] (0,-0.05) -- (0,1.15);
        \node[below, font=\tiny] at (0,-0.05) {$0$};
        % frecuencia de Nyquist, rotulada abajo y en segunda fila para no chocar con wr
        \draw[nyq] (3.5,-0.28) -- (3.5,1.24);
        \draw[nyq] (-3.5,-0.28) -- (-3.5,1.24);
        \node[gray!55!black, font=\tiny, anchor=north] at (3.5,-0.30) {$\tfrac{\omega_s}{2}$};
        % marcas del eje
        \draw[thin] (1.8,0.04) -- (1.8,-0.04) node[below, font=\tiny] {$\omega_M$};
        \draw[thin] (3.0,0.04) -- (3.0,-0.04) node[below, font=\tiny] {$\omega_r$};
        % el rótulo va en una segunda fila, con guía, para no confundirlo con ws/2
        \draw[thin] (4.0,0.04) -- (4.0,-0.04);
        \node[anchor=west, font=\tiny] at (4.6,-0.36) {$\omega_s - \omega_r$};
        \draw[thin, -{Latex[length=3pt]}] (4.5,-0.34) -- (4.06,-0.10);
        \draw[thin] (7.0,0.04) -- (7.0,-0.04) node[below, font=\tiny] {$\omega_s$};
        % regla superior: banda útil, transición y franja libre, una tras otra
        \draw[med] (-1.8,1.30) -- (1.8,1.30);
        \node[gray!65, anchor=south, font=\tiny, fill=white, inner sep=1pt] at (0,1.35) {banda útil};
        \draw[med] (1.8,1.30) -- (3.0,1.30);
        \node[gray!65, anchor=south, font=\tiny, fill=white, inner sep=1pt] at (2.4,1.52) {transición};
        \draw[med] (3.0,1.30) -- (4.0,1.30);
        \node[gray!65, anchor=south, font=\tiny, fill=white, inner sep=1pt] at (3.5,1.35) {libre};
        \node[utnblue!70, anchor=south, font=\tiny] at (-7.0,1.06) {vecina};
        \draw[utnblue!70, thin, -{Latex[length=3pt]}] (-7.0,1.04) -- (-7.0,0.30);
    \end{scope}
\end{tikzpicture}
\end{document}
